\documentclass[10pt]{article}
\providecommand{\FIGDIR}{fig}
\newcommand{\DIGESTDATE}{25 August 2026}
% ===================== PM Research Watch - shared preamble =====================
\usepackage[T1]{fontenc}
\usepackage[utf8]{inputenc}
\usepackage{amsmath,amssymb}
\usepackage{tgpagella}
\usepackage{tgheros}
\renewcommand{\sfdefault}{qhv}
\renewcommand{\ttdefault}{lmtt}
\usepackage{microtype}
\usepackage[english]{babel}

\usepackage[a4paper,top=17mm,bottom=17mm,left=17mm,right=17mm,
            headsep=6mm,footskip=10mm]{geometry}
\usepackage{graphicx}
\usepackage{xcolor}
\usepackage{colortbl}
\usepackage{tikz}
\usetikzlibrary{calc}
\usepackage[most]{tcolorbox}
\usepackage{titlesec}
\usepackage{fancyhdr}
\usepackage{booktabs}
\usepackage{array}
\usepackage{longtable}
\usepackage{enumitem}
\usepackage{multicol}
\usepackage{ragged2e}
\usepackage[hidelinks]{hyperref}

% ------------------------------------------------------------------ palette
\definecolor{Ink}    {HTML}{14262E}
\definecolor{Deep}   {HTML}{0F4C5C}
\definecolor{Teal}   {HTML}{1F7A8C}
\definecolor{Sky}    {HTML}{4FA3C4}
\definecolor{Sage}   {HTML}{6FA287}
\definecolor{Amber}  {HTML}{E3A008}
\definecolor{Clay}   {HTML}{B5643C}
\definecolor{Coral}  {HTML}{D1495B}
\definecolor{Violet} {HTML}{6A5B9E}
\definecolor{Slate}  {HTML}{4A6572}
\definecolor{Paper}  {HTML}{FBF9F5}
\definecolor{Mist}   {HTML}{EEEAE2}
\definecolor{Rule}   {HTML}{D9D3C9}

\pagecolor{Paper}
\color{Ink}

\hypersetup{colorlinks=true, linkcolor=Deep, urlcolor=Teal, citecolor=Teal,
            pdftitle={PM Research Watch}, pdfauthor={Automated daily digest}}

% ------------------------------------------------------------------ headings
\titleformat{\section}
  {\sffamily\bfseries\large\color{Deep}}{}{0pt}
  {}[{\vspace{1.2mm}\color{Amber}\titlerule[1.1pt]}]
\titlespacing*{\section}{0pt}{5.5mm}{2.8mm}

\titleformat{\subsection}
  {\sffamily\bfseries\normalsize\color{Teal}}{}{0pt}{}
\titlespacing*{\subsection}{0pt}{3.6mm}{1.4mm}

\setlist[itemize]{leftmargin=4.2mm, itemsep=1.1mm, topsep=1.3mm, parsep=0pt,
                  label=\textcolor{Teal}{\footnotesize$\blacktriangleright$}}

% ------------------------------------------------------------------ header/footer
\pagestyle{fancy}
\fancyhf{}
\renewcommand{\headrulewidth}{0pt}
\renewcommand{\footrulewidth}{0pt}
\fancyhead[L]{\sffamily\scriptsize\color{Slate}PM RESEARCH WATCH\;\textcolor{Rule}{\textbar}\;\DIGESTDATE}
\fancyhead[R]{\sffamily\scriptsize\color{Slate}Particulate matter monitoring \& health effects}
\fancyfoot[C]{\sffamily\scriptsize\color{Slate}\thepage}

% ------------------------------------------------------------------ components
% metric tile
\newcommand{\metric}[3]{%
  \begin{tcolorbox}[enhanced, colback=white, colframe=Rule, boxrule=0.5pt,
    arc=1.6mm, left=2.2mm, right=2.2mm, top=1.8mm, bottom=1.8mm,
    borderline west={1.7mm}{0pt}{#3}, nobeforeafter, valign=center,
    height=17mm, width=\linewidth]
    \RaggedRight
    {\sffamily\bfseries\fontsize{17}{17}\selectfont\color{#3}#1}\\[0.5mm]
    {\sffamily\fontsize{6.5}{7.4}\selectfont\color{Slate}#2}
  \end{tcolorbox}}

% pull-quote / key insight box
\newtcolorbox{keybox}[1][]{enhanced, colback=Mist, colframe=Mist,
  arc=1.6mm, left=3.2mm, right=3.2mm, top=2.4mm, bottom=2.4mm,
  borderline west={1.7mm}{0pt}{Amber}, #1}

% subtopic band
\newcommand{\band}[2]{%
  \vspace{3.2mm}\noindent
  \begin{tcolorbox}[enhanced, colback=#2, colframe=#2, arc=1.2mm,
    left=2.6mm, right=2.6mm, top=1.1mm, bottom=1.1mm, boxrule=0pt,
    width=\linewidth, nobeforeafter]
    {\sffamily\bfseries\small\color{white}#1}
  \end{tcolorbox}\vspace{1.6mm}}

% one paper entry: {short cite}{title}{journal}{doi}{body}
\newcommand{\paperentry}[5]{%
  \par\vspace{1.9mm}\noindent
  {\sffamily\bfseries\footnotesize\color{Ink}#1}\;%
  {\sffamily\fontsize{7}{8}\selectfont\color{Slate}#3}\\[0.4mm]
  {\footnotesize\itshape\color{Deep}#2}\\[0.6mm]
  {\footnotesize #4}\\[0.3mm]
  {\fontsize{6.6}{7.6}\selectfont\color{Slate}\texttt{doi:}\,\href{https://doi.org/#5}{\texttt{\detokenize{#5}}}}
  \par}

% caption style
\newcommand{\figcap}[1]{\par\vspace{1.0mm}
  {\fontsize{7.2}{8.6}\selectfont\color{Slate}\textsf{\textbf{Fig.}}\; #1\par}\vspace{1.5mm}}

\providecommand{\permil}{\textperthousand}

\setlength{\parindent}{0pt}
\setlength{\parskip}{0pt}
\renewcommand{\arraystretch}{1.18}


\begin{document}

% ==================================================================== MASTHEAD
\thispagestyle{empty}
\begin{tikzpicture}[remember picture, overlay]
  \fill[Deep] (current page.north west) rectangle
      ($(current page.north east)+(0,-67mm)$);
  \fill[Amber] ($(current page.north west)+(0,-67mm)$) rectangle
      ($(current page.north east)+(0,-68.4mm)$);
  \foreach \x/\y/\r in {22/12/0.9, 41/26/1.5, 63/9/0.7, 88/31/1.1, 112/17/1.9,
                        138/29/0.8, 159/13/1.3, 178/54/1.0, 196/21/1.6,
                        30/44/0.6, 74/57/1.2, 125/48/0.9, 168/6/0.8, 191/59/0.7,
                        52/18/1.0, 100/60/0.6, 148/20/1.1, 66/49/1.4, 205/45/0.8}
    \fill[Sky, opacity=0.30] ($(current page.north west)+(\x mm,-\y mm)$) circle (\r mm);
\end{tikzpicture}

\vspace*{4mm}
{\sffamily\fontsize{7.6}{9}\selectfont\color{Sky}%
 AUTOMATED DAILY LITERATURE DIGEST\;\textbullet\;PhD RESEARCH WATCH\par}
\vspace{2.2mm}
{\sffamily\bfseries\fontsize{27}{29}\selectfont\color{white}%
 Particulate Matter\par}
{\sffamily\fontsize{27}{29}\selectfont\color{Sky}%
 Monitoring \& Health Effects\par}
\vspace{5.0mm}
{\sffamily\fontsize{8.4}{10}\selectfont\color{white}%
 \mbox{Issue 2026-08-25}\;\;\textcolor{Amber}{\textbar}\;\;\mbox{Entry window 25 Aug 2026}\;\;
 \textcolor{Amber}{\textbar}\;\;\mbox{14 in scope, 51 screened out}\par}
\vspace{18mm}

% -------------------------------------------------------------- metric strip
\noindent
\begin{minipage}[t]{0.190\linewidth}\metric{14}{RECORDS IN SCOPE\\(51 SCREENED OUT)}{Deep}\end{minipage}\hfill
\begin{minipage}[t]{0.190\linewidth}\metric{7}{OF 10 CLUSTERS\\OCCUPIED}{Teal}\end{minipage}\hfill
\begin{minipage}[t]{0.190\linewidth}\metric{5}{PREPRINT-STAGE\\RECORDS (36\%)}{Sky}\end{minipage}\hfill
\begin{minipage}[t]{0.190\linewidth}\metric{4}{INTERVAL-BEARING\\ESTIMATES (1 STUDY)}{Amber}\end{minipage}\hfill
\begin{minipage}[t]{0.190\linewidth}\metric{1}{ABSTRACT-FREE\\DEPOSIT (7\%)}{Clay}\end{minipage}

\vspace{6mm}

% -------------------------------------------------------------- signal of day
\section*{\sffamily\bfseries\color{Deep}Signal of the day}
\vspace{-1.5mm}{\color{Amber}\rule{\linewidth}{1.1pt}}\vspace{2.5mm}

\begin{keybox}
{\sffamily\bfseries\small\color{Deep}A full year of riverbank co-location shows that the
binding constraint on plume-resolved low-cost sensing is not correlation with the
reference --- it is noise floor and response time.}\\[1.4mm]
\footnotesize
Eger et al.\ (\textit{Aerosol Air Qual Res}) co-located a panel of low- and mid-cost
sensors against reference instrumentation on the Rhine for one year, screening them
specifically for the ability to resolve individual inland-ship exhaust plumes.
Correlation was never the problem: \textbf{MA350 black carbon reached $R^2=0.95$
(slope 0.77), Partector 2 Pro ultrafine particle number $R^2=0.85$, K30 FR CO$_2$
$R^2=0.98$, and the Alphasense NO and NO$_2$ cells $R^2=0.99$ and $0.81$}. What limited
the study was that the sensors' instrumental noise and the temperature dependence of
the electrochemical cells reduced the number of ship peaks that could be analysed at
all. This is the distinction that urban-network validation studies systematically
obscure: hourly-averaged agreement with a reference monitor says almost nothing about
whether a sensor can resolve a transient. Any deployment aimed at source attribution
rather than compliance-style mean concentration should be specified on
\emph{noise floor and $t_{90}$}, not on $R^2$.
\end{keybox}

\vspace{3mm}

% -------------------------------------------------------------- exec brief
\section*{\sffamily\bfseries\color{Deep}What changed today}
\vspace{-1.5mm}{\color{Amber}\rule{\linewidth}{1.1pt}}\vspace{2.5mm}

\footnotesize
\begin{itemize}\itemsep2pt
\item \textbf{Five of fourteen records are preprints}, the highest preprint share this
watch has carried. All five arrive through the Europe PMC leg (three medRxiv, two
Research Square); none has been peer reviewed. They are tiered and summarised on the same terms as the
journal records but the tier reflects design, not review status --- read the
Mandalapu, Walker, Pathak and Fahim entries accordingly.

\item \textbf{Two independent preprints argue that the exposure \emph{metric}, not the
exposure, is what determines the disparity you observe.} Mandalapu et al.\ find that
13.21\,\% of residents of predominantly Black tracts live in joint heat-plus-PM$_{2.5}$
hotspots against 3.33\,\% in predominantly White tracts, with \textbf{OR 3.97
(95\,\% CI 2.79--5.66)} for HOLC grade D versus grade A after adjustment. Walker et al.,
measuring directly at 176 Rhode Island sites, find that the social gradient sits in the
acoustic environment and night-time light --- \textbf{short-term PM$_{2.5}$ was roughly an
order of magnitude weaker} and showed no reliable gradient. Both are correct; they differ
in averaging time and in whether the metric is modelled-annual or measured-instantaneous.
Composite environmental-justice indices that mix the two are not measuring one thing.

\item \textbf{A foundation-model result with a directly usable negative finding.}
Fahim \& Karim benchmark zero-shot time-series foundation models against tuned
specialists across 29 cities. On PM$_{2.5}$ the two are statistically indistinguishable
(MASE 0.662 vs 0.692), transfer learning never beats zero-shot at any budget, and under
frequent retraining the untrained model uses roughly ten times less measured energy.
More useful still: the specialist's apparent temperature advantage (0.533 vs 0.792)
collapses to a non-significant 0.745 once its covariates are restricted to values
knowable at forecast time --- a \emph{perfect-foresight artefact} that will be present in
any PM forecasting benchmark that feeds future meteorology to one arm and not the other.

\item \textbf{mtDNA copy number is not yet a usable PM exposure biomarker.} Pathak et al.\
pool 24 studies and 13,092 participants and find PM associated with a
\textbf{$-4.90\,\%$ (95\,\% CI $-7.97$ to $-1.82$) change in percentage mtDNA-CN} but a
\textbf{$+0.55$ (0.05--1.04) change in absolute mtDNA-CN}. The directions disagree
because the metrics disagree; that is a measurement-harmonisation failure in the
underlying literature, not a biological finding.

\item \textbf{Indoor combustion, interaction terms, and ventilation.} Chen et al.\
monitored 24\,h indoor PM$_{2.5}$/PM$_{10}$ and gases in 60 Taiwanese households and
report that PM$_{2.5}$$\times$TVOC and PM$_{2.5}$$\times$CO interactions --- not
PM$_{2.5}$ alone --- carried the adverse effects on myocardial contractility and vascular
elasticity, with elevated CO$_2$ acting as an independent stressor in
non-incense-burning homes. A PM-only indoor monitor cannot see this.

\item \textbf{One abstract-free deposit.} Vivanco et al.\ (\textit{Environ Pollut}) on
air-quality-model responsiveness to real-world emission changes arrived with title and
DOI but no abstract, and is carried at tier C on the title alone. The Elsevier
abstract backlog now spans four weeks.
\end{itemize}

\clearpage

% ==================================================================== ANALYTICS
\section*{\sffamily\bfseries\color{Deep}Corpus at a glance}
\vspace{-1.5mm}{\color{Amber}\rule{\linewidth}{1.1pt}}\vspace{3mm}

\noindent\includegraphics[width=\linewidth]{\FIGDIR/f1_subtopics.png}
\figcap{Records per subtopic cluster. Seven of ten clusters are occupied.
\textit{Sensing, forecasting \& instrumentation}, \textit{Exposure assessment \&
modelling} and \textit{Mechanistic toxicology} tie at 3 records each;
every mechanistic record is a synthesis rather than a primary experiment.
Assignment is single-label by dominant contribution.}

\vspace{2mm}

\vspace{2mm}
\noindent\includegraphics[width=\linewidth]{\FIGDIR/f2_design_endpoint.png}
\figcap{Left --- study architecture as binned by the coarse grouper.
\textit{Review / synthesis} leads at 4, with \textit{Measurement campaign} and
\textit{Observational --- acute} tied at 3; \textit{Modelling / inventory} 2, and one
each for \textit{Ecological} and \textit{Observational --- cross-sectional}. Note that
Eger et al.\ --- a year-long instrument co-location, not an epidemiological field study
--- is binned as a measurement campaign, the closest available group rather than an
accurate description. Right --- health endpoint. \textbf{Six of the fourteen records
carry no human health endpoint at all}: the three sensing/forecasting records and the
three exposure-assessment records report concentration, hazard quotient or forecast
skill only. The remaining eight endpoints are singletons --- no two records in this
window share an outcome.}


\vspace{2mm}
\noindent\includegraphics[width=\linewidth]{\FIGDIR/f3_metric_geography.png}
\figcap{Left --- particle metric. \textit{PM$_{2.5}$ only} 5 and
\textit{PM$_{2.5}$ + PM$_{10}$ jointly} 4 dominate; \textit{PM unspeciated / emissions}
3; \textit{Bioaerosol / coarse} and \textit{Ultrafine / nanoscale} 1 each. The grouper
resolved Eger et al.\ (black carbon plus ultrafine particle \emph{number}) and Yang
et al.\ (culturable bioaerosol counts) into their own bins rather than folding them into
a mass bin. Right --- geographic provenance. \textbf{China alone supplies 5 of 14
records}; multi-region syntheses 3, North America 2, Europe 2, and one each from East
Asia outside China (Taiwan) and Southeast Asia (Thailand). Nothing from Africa, South
America or South Asia.}


\vspace{2mm}
\noindent\includegraphics[width=\linewidth]{\FIGDIR/f4_heatmap.png}
\figcap{Subtopic $\times$ study-architecture cross-tabulation. The densest cell is
\textit{Mechanistic toxicology} $\times$ \textit{Review / synthesis} (3) --- every
mechanistic record in the window is a synthesis, and not one is a primary experiment.
The \textit{Sensing} row is split one-and-two across instrument evaluation and
modelling, and the \textit{Exposure assessment} row contains no health-endpoint
linkage whatsoever.}

\vspace{2mm}
\noindent\includegraphics[width=\linewidth]{\FIGDIR/f5_forest.png}
\figcap{Every interval-bearing estimate in the window, and \textbf{all four come from a
single unreviewed preprint} (Mandalapu et al.). The top estimate is an adjusted odds
ratio for joint-hotspot residence; the three below it are \emph{conditional exceedance
ratios} --- observed joint occurrence against within-group statistical independence ---
and are therefore not on the same scale as the odds ratio and not comparable to it.
The Black and Hispanic intervals both include 1; the authors attribute this to limited
power from fewer contributing metropolitan areas rather than to an absence of
dependence. Read the plot for precision and direction only.}


\vspace{2mm}
\noindent\includegraphics[width=\linewidth]{\FIGDIR/f6_lifecourse.png}
\figcap{Life-course windows addressed, counted only where a record's study population is
explicitly anchored to a stage. Three records qualify, all working-age adults. Ten
records study no human population at all or pool across ages without stratification;
the in-utero, paediatric, adolescent and 65+ windows are empty for the second
consecutive issue.}


\clearpage

% ==================================================================== DIGEST
\section*{\sffamily\bfseries\color{Deep}Paper-level digest}
\vspace{-1.5mm}{\color{Amber}\rule{\linewidth}{1.1pt}}\vspace{1mm}

{\fontsize{7.4}{9}\selectfont\color{Slate}Ordered by cluster. Each entry gives the
methodological core and the finding that matters, not an abstract paraphrase.
Preprints are labelled as such in the journal field.\par}

\band{Sensing, forecasting \& instrumentation\hfill 3 records}{Amber}

\paperentry{Eger et al.\ 2026}
{Compact low- and mid-cost sensors for inland-ship exhaust emissions}
{Aerosol Air Qual Res}
{One year of onshore co-location on the Rhine against reference instruments, screening
sensors on noise floor, correlation and response time --- criteria chosen because ship
plumes are transients, not the slowly-varying urban background most validation studies
target. Qualifying devices: \textbf{MA350 black carbon ($R^2=0.95$, slope 0.77),
Partector 2 Pro ultrafine particle number ($R^2=0.85$), K30 FR CO$_2$ ($R^2=0.98$),
Alphasense NO ($R^2=0.99$) and NO$_2$ ($R^2=0.81$)}. The limitation the authors state
plainly is the useful part: sensor noise and the temperature dependence of the
electrochemical cells --- which they describe as difficult to correct --- cut the number
of analysable ship peaks. No PM mass sensor cleared the bar, which is itself a result.}
{10.1007/s44408-026-00163-w}

\paperentry{Fahim \& Karim 2026}
{Cost-aware benchmark of time-series foundation models for PM$_{2.5}$ forecasting}
{Res Sq (preprint)}
{Twenty-nine cities, two domains (hourly PM$_{2.5}$, 2\,m temperature), pre-specified
protocols with measured energy and a cost-adjusted decision rule, comparing zero-shot
foundation models, transfer learning and tuned specialists. \textbf{On PM$_{2.5}$ the
zero-shot model and the specialist are statistically indistinguishable (MASE 0.662 vs
0.692)}; transfer learning never significantly beats zero-shot at any budget; frequent
retraining costs the specialist roughly ten times the measured energy. The
methodological contribution is the demonstration that the specialist's temperature
advantage is a \emph{perfect-foresight artefact} of feeding it covariates unknowable at
forecast time. Unreviewed, and 29 cities is a small sample for a claim about
generalisation --- but the artefact critique applies to published benchmarks too.}
{10.21203/rs.3.rs-10515545/v1}

\paperentry{Vivanco et al.\ 2026}
{Responsiveness of an air quality model to real-world emission changes}
{Environ Pollut}
{\textbf{Metadata-only deposit} --- title and DOI without an abstract, so the design,
model, domain and emission perturbations are all unknown at the time of writing. The
question is a live one (whether CTMs used for planning actually reproduce the
concentration response to observed emission changes rather than only to scenario
deltas) and the record is carried at tier C so it is not lost. Flagged for
abstract retry.}
{10.1016/j.envpol.2026.129031}

\band{Exposure assessment \& modelling\hfill 3 records}{Deep}

\paperentry{Mandalapu et al.\ 2026}
{Joint heat and PM$_{2.5}$ co-exposure across 48 US metropolitan areas}
{medRxiv (preprint)}
{42,304 census tracts, 48 large CBSAs, ${\sim}174.6$ million residents, summers
2015--2020. The framework combines a composite co-exposure index, a joint exceedance
indicator, a conditional exceedance ratio against within-group independence, and an
upper-tail dependence parameter estimated both non-parametrically
(Cap\'era\`a--Foug\`eres--Genest) and with a Gumbel copula, with BCa intervals from a
5,000-replicate block bootstrap. \textbf{13.21\,\% of residents of predominantly Black
tracts versus 3.33\,\% of predominantly White tracts live in joint hotspots;
HOLC grade D vs grade A gives OR 3.97 (2.79--5.66)} after adjustment for contemporary
racial composition, poverty, renter occupancy and pre-1960 housing. Magnitudes attenuate
under near-surface air temperature rather than surface temperature. Unreviewed;
the tail-dependence machinery is the transferable contribution.}
{10.64898/2026.08.20.26360970}

\paperentry{Walker et al.\ 2026}
{Directly measured noise, PM$_{2.5}$ and night light across Rhode Island}
{medRxiv (preprint)}
{176 sites, 704 site-sessions over roughly a year, spanning the Providence--Central
Falls--Pawtucket urban core and four rural municipalities, with A- and C-weighted noise,
PM$_{2.5}$, night-time illuminance and relative humidity carried through an identical
model as a negative control. A-weighted noise rose with disadvantage on six of eight
area-level markers (\textbf{1.3--1.8\,dBA per SD; 6.6\,dBA between least and most
racially diverse neighbourhoods}) and C-weighted levels followed the same gradient,
exceeding their A-weighted counterparts at block-group scale for renter occupancy and
vehicle absence. \textbf{Short-term PM$_{2.5}$ was roughly an order of magnitude
weaker}; humidity showed nothing. The negative-control design is what makes the PM
null interpretable rather than merely underpowered. Short-term measurement cannot
speak to annual-average disparity, and the authors do not claim it does.}
{10.64898/2026.08.20.26360933}

\paperentry{Mueangsai et al.\ 2026}
{PM$_{2.5}$ health risk quotient near stone crushing plants, Chonburi}
{Res Sq (preprint)}
{Community cross-sectional study, 368 residents by stratified random sampling, with
PM$_{2.5}$ from two fixed stations near a crushing plant averaged and pushed through the
USEPA health risk assessment to a hazard quotient. The design's weakness is structural
and fatal to individual-level inference: \textbf{two fixed monitors supply a single
exposure value for every participant}, so all between-person exposure contrast is
assumed away and the HQ varies only through the questionnaire covariates. Carried at
tier C. The setting --- crushing plants generate coarse, mineral-dominated PM that
optical sensors size poorly --- is the reason it is here at all.}
{10.21203/rs.3.rs-10160540/v1}

\clearpage

\band{Mechanistic toxicology\hfill 3 records}{Violet}

\paperentry{Pathak et al.\ 2026}
{Airborne PM and mitochondrial DNA copy number: systematic review and meta-analysis}
{medRxiv (preprint)}
{PRISMA, PROSPERO-registered, JBI risk-of-bias, random-effects pooling; 1,224 records
screened to 24 studies and \textbf{13,092 participants}, 12 reporting percentage change
and 12 absolute values. \textbf{Percentage mtDNA-CN fell ($-4.90$; 95\,\% CI $-7.97$ to
$-1.82$; $p=0.002$) while absolute mtDNA-CN rose ($+0.55$; 0.05--1.04; $p=0.030$).}
The two directions are not reconcilable as biology; they are an artefact of pooling
across incompatible normalisation schemes, and the honest reading is that the
underlying literature has not standardised its outcome. Useful as evidence that
mtDNA-CN is \emph{not yet} deployable as a PM exposure biomarker, which is the opposite
of how it is usually cited.}
{10.64898/2026.08.20.26360559}

\paperentry{Forest 2026}
{Molecular mechanisms of sand and dust storm toxicity}
{Sci Total Environ}
{Review of in vitro and in vivo evidence on storm-derived PM$_{10}$/PM$_{2.5}$, whose
composition --- quartz and clay minerals plus adsorbed bacteria, viruses and heavy
metals acquired in transport --- differs enough from urban PM that the author argues
current regulation, which treats storm episodes as generic PM, misses hazard-specific
risk. Consistent findings are cytotoxicity, oxidative stress, pro-inflammatory response
and exacerbation of bacterial and viral infection. The stated gaps are the substantive
content: \textbf{most in vivo work uses Asian sand dust while Saharan dust dominates
global emissions}, and non-respiratory target organs are barely studied. Directly
relevant to any network operating in a dust-influenced region, where the mass a sensor
reports and the hazard it implies decouple.}
{10.1016/j.scitotenv.2026.182233}

\paperentry{Li \& Li 2026}
{GPER-mediated pulmonary carcinogenesis by fine particulate matter}
{Zhongguo Fei Ai Za Zhi}
{Chinese-language narrative review proposing that estrogen-like organic components and
metal ions in PM$_{2.5}$ activate the G protein-coupled estrogen receptor, which is
overexpressed in lung tumour tissue, driving MAPK/ERK and downstream signalling. No new
data and no quantitative synthesis; the receptor-level specificity is the only reason
to prefer it to the generic PM-and-lung-cancer literature. Tier C.}
{10.3779/j.issn.1009-3419.2026.101.13}

\band{Cardiovascular \& metabolic\hfill 1 record}{Teal}

\paperentry{Chen et al.\ 2026}
{Residential incense burning, indoor pollutants and cardiovascular haemodynamics}
{Indoor Air}
{Seventy-three adults from the YOung TAiwanese cohort across 60 households, with
continuous 24\,h indoor PM$_{2.5}$, PM$_{10}$, CO, CO$_2$, TVOC and formaldehyde
alongside non-invasive haemodynamic measurement, analysed by mixed-effects models.
Incense-burning households showed higher pollutant concentrations and worse
haemodynamics --- higher blood pressure and cardiac workload, lower vascular compliance.
\textbf{The finding that matters is interactive: PM$_{2.5}$$\times$TVOC and
PM$_{2.5}$$\times$CO carried the effects on myocardial contractility and vascular
elasticity}, and in non-incense homes elevated CO$_2$ acted as an independent stressor.
No effect sizes are given in the abstract. $n=73$ with household clustering is thin for
interaction terms, and residual confounding by household behaviour is not addressable
here. Implication for instrumentation: a PM-only indoor node is blind to the
combination that did the work.}
{10.1155/ina/8778576}

\band{Neuro / mental health\hfill 1 record}{Sky}

\paperentry{Ma et al.\ 2026}
{Day-on-day increases in air pollution and depression morbidity, 239 Chinese cities}
{Ecotoxicol Environ Saf}
{Time-stratified case-crossover on \textbf{330,610 hospital admissions and 3,726,390
outpatient visits} for depression, 239 prefecture-level-or-above cities, 2013--2017.
The exposure metric is the novelty: air pollution increase between neighbouring days
(APIN), the current day's absolute concentration minus the previous day's, rather than
level. Per IQR increase in APIN at lag 0--3, admissions rose between 0.59\,\%
(0.06--1.13) for 8\,h O$_3$ and 1.66\,\% (1.09--2.22) for NO$_2$, with outpatient visits
0.42--0.85\,\%; associations were robust to adjustment for absolute concentrations
except for ozone. \textbf{PM$_{2.5}$- and PM$_{10}$-specific estimates are not reported
in the abstract}, only the range endpoints, so no PM estimate is carried into the forest
plot. The rate-of-change framing is worth borrowing regardless of the endpoint.}
{10.1016/j.ecoenv.2026.120707}

\clearpage

\band{Respiratory \& allergic\hfill 2 records}{Sage}

\paperentry{China Lung Oncology Group 2026}
{Regional lung cancer related to coal pollution, Yunnan--Guizhou Plateau}
{Zhongguo Fei Ai Za Zhi}
{Multi-sectional, multi-omics regional evidence synthesis proposing ``China Wumeng
Mountain lung cancer'' as a distinct entity driven by long-term household coal
combustion, with the argument resting on differences in coal source and mineral
composition rather than on total exposure. Clinical and molecular characterisation is
paired with a proposed prevention pathway. The construct is at risk of circularity ---
defining a regional cancer type by its presumed cause and then citing regional
incidence as evidence --- and the abstract gives no effect estimates. Tier B on the
strength of the source-composition argument, which is the same argument the wildfire and
dust literatures are making independently.}
{10.3779/j.issn.1009-3419.2026.102.16}

\paperentry{Yang et al.\ 2026a}
{Spatiotemporal risk of environmental drivers of respiratory disease, Jiangsu}
{GeoHealth}
{Distributed lag non-linear models coupled to geographically and temporally weighted
regression on respiratory hospitalisations across Jiangsu Province, 2020--2023, a
population containing over 27 million children under 15 and adults over 65. The DLNM
handles the lag structure and the GTWR lets the exposure--response coefficient vary in
space and time rather than assuming one province-wide slope --- which is the right
question for a province with this much industrial heterogeneity. The 2020--2023 window
overlaps pandemic-era mobility and healthcare-seeking disruption; the abstract does not
say how that is handled, and it is the main threat to the temporal component.}
{10.1029/2026gh001881}

\band{Occupational \& indoor\hfill 1 record}{Clay}

\paperentry{Yang et al.\ 2026b}
{Bioaerosols at oil exploitation sites and mask filtration efficiency}
{Environ Pollut}
{Seasonal sampling in a Beijing--Tianjin--Hebei oil city carrying both heavy PM and VOC
loads, characterising culturable and total bioaerosol distribution, sources and
inhalation risk, plus mask filtration efficiency. \textbf{Autumn levels significantly
exceeded summer while microbial activity moved the other way} --- a dissociation between
count and viability that matters for anyone interpreting a fluorescence channel as a
proxy for biological hazard. Culture-based enumeration undercounts viable-but-non-culturable
organisms by a large and site-dependent factor, which is the standing limitation of this
entire measurement class. Tier C.}
{10.1016/j.envpol.2026.129020}

\clearpage

% ==================================================================== TABLE
\section*{\sffamily\bfseries\color{Deep}Corpus register}
\vspace{-1.5mm}{\color{Amber}\rule{\linewidth}{1.1pt}}\vspace{2.5mm}

{\fontsize{6.7}{8.4}\selectfont
\setlength{\tabcolsep}{2.4pt}
\begin{longtable}{@{}>{\RaggedRight}p{25mm}>{\RaggedRight}p{41mm}>{\RaggedRight}p{27mm}%
>{\RaggedRight}p{23mm}>{\RaggedRight}p{18mm}>{\RaggedRight}p{15mm}@{}}
\toprule
\textsf{\textbf{First author}} & \textsf{\textbf{Focus}} & \textsf{\textbf{Journal}} &
\textsf{\textbf{Design}} & \textsf{\textbf{Metric}} & \textsf{\textbf{Region}}\\
\midrule
\endfirsthead
\toprule
\textsf{\textbf{First author}} & \textsf{\textbf{Focus}} & \textsf{\textbf{Journal}} &
\textsf{\textbf{Design}} & \textsf{\textbf{Metric}} & \textsf{\textbf{Region}}\\
\midrule
\endhead
\multicolumn{6}{@{}l@{}}{\cellcolor{Mist}\textsf{\textbf{\textcolor{Deep}{Neuro / mental health}}}}\\[0.6mm]
\href{https://doi.org/10.1001/jamanetworkopen.2026.25015}{Dales} & Wildfire-sourced PM$_{2.5}$ and ED visits for migraine and other primary headache & \textit{JAMA Netw Open} & Case-crossover & PM$_{2.5}$ (wildfire) & Canada\\
\href{https://doi.org/10.1093/ageing/afag215}{Glover} & Umbrella meta-review of urban environmental exposures and cognitive health & \textit{Age Ageing} & Umbrella review / meta-analysis & PM$_{2.5}$ & Global\\
\href{https://doi.org/10.1016/j.nsa.2026.107019}{Goossens} & In utero air pollution exposure and pre-/postnatal brain development & \textit{Neurosci Appl} & Narrative review & PM$_{2.5}$ / PM$_{10}$ / PM$_{\mathrm{coarse}}$ & Global\\
\href{https://doi.org/10.1016/j.isci.2026.116617}{Ma} & PM$_{2.5}$-induced neurotoxicity and depression: exposure risks and mechanisms & \textit{iScience} & Narrative review & PM$_{2.5}$ + components & Global\\
\href{https://doi.org/10.3760/cma.j.cn112150-20250804-00752}{Pan} & Maternal PM exposure and offspring neurodevelopment: effects and mechanisms & \textit{Zhonghua Yu Fang Yi Xue Za Zhi} & Narrative review & PM (bulk) & China\\
\addlinespace[0.9mm]
\multicolumn{6}{@{}l@{}}{\cellcolor{Mist}\textsf{\textbf{\textcolor{Teal}{Cardiovascular \& metabolic}}}}\\[0.6mm]
\href{https://doi.org/10.3760/cma.j.cn112338-20251120-00843}{Huang} & PM$_{2.5}$ and dyslipidaemia in rural Chinese children; dietary effect modification & \textit{Zhonghua Liu Xing Bing Xue Za Zhi} & Cross-sectional (multistage) & PM$_{2.5}$ & China\\
\href{https://doi.org/10.32394/pe/218966}{Kazmierska} & Air pollution and cardiovascular disease: review of current data & \textit{Przegl Epidemiol} & Narrative review & PM$_{2.5}$ / PM$_{10}$ & Global\\
\href{https://doi.org/10.1007/s00125-026-06814-2}{Xi} & Air pollution, built environment and temperature vs. prediabetes phenotypes & \textit{Diabetologia} & Prospective cohort & PM$_{2.5}$ & Germany\\
\href{https://doi.org/10.1161/JAHA.125.046917}{Yang} & Short-term PM$_{2.5}$ and hospitalisation/ED use in atrial fibrillation & \textit{J Am Heart Assoc} & Case-crossover & PM$_{2.5}$ (wildfire-influenced) & USA\\
\addlinespace[0.9mm]
\multicolumn{6}{@{}l@{}}{\cellcolor{Mist}\textsf{\textbf{\textcolor{Sky}{Reproductive \& developmental}}}}\\[0.6mm]
\href{https://doi.org/10.1016/j.jes.2025.10.043}{Chen} & Prenatal air pollution, fetal growth trajectory phenotypes and child neurodevelopment & \textit{J Environ Sci (China)} & Prospective cohort & PM$_{2.5}$ / PM$_{10}$ & China\\
\href{https://doi.org/10.1016/j.ecoenv.2026.120566}{Huang} & Joint prenatal bisphenol + air-pollutant mixture exposure and preschool IQ & \textit{Ecotoxicol Environ Saf} & Prospective cohort (mixtures) & PM$_{2.5}$ / PM$_{10}$ & China\\
\href{https://doi.org/10.1007/s40572-026-00552-8}{Peterson} & Ambient air pollution, PAHs and pubertal development: critical review & \textit{Curr Environ Health Rep} & Critical review & PM$_{2.5}$ components & Global\\
\href{https://doi.org/10.1016/j.envres.2026.125313}{Xie} & Long-term PM$_{2.5}$ and timing of pubertal development in boys and girls & \textit{Environ Res} & Prospective cohort & PM$_{2.5}$ & China\\
\addlinespace[0.9mm]
\multicolumn{6}{@{}l@{}}{\cellcolor{Mist}\textsf{\textbf{\textcolor{Sage}{Respiratory \& allergic}}}}\\[0.6mm]
\href{https://doi.org/10.1007/s00484-026-03289-z}{Bayar Muluk} & Climate change and otolaryngology: airway disease impacts and emission reduction & \textit{Int J Biometeorol} & Narrative review & PM$_{2.5}$ & Global\\
\href{https://doi.org/10.4103/mgr.MEDGASRES-D-25-00164}{Liao} & Short-term air pollution, respiratory symptoms and lung function in asthma & \textit{Med Gas Res} & Panel study & PM$_{2.5}$ / PM$_{10}$ & China\\
\href{https://doi.org/10.3390/epidemiologia7040097}{Sanchez de Prada} & Ambient PM and asthma-coded admissions in children 0-3 y, Spain & \textit{Epidemiologia} & Ecological panel & PM$_{2.5}$ / PM$_{10}$ & Spain\\
\href{https://doi.org/10.1016/j.jhazmat.2026.143076}{Yang} & Allergenic vs. toxic risk of airborne seed hairs compared with ambient PM$_{2.5}$ & \textit{J Hazard Mater} & Field measurement & PM$_{2.5}$ + bioaerosol & China\\
\addlinespace[0.9mm]
\multicolumn{6}{@{}l@{}}{\cellcolor{Mist}\textsf{\textbf{\textcolor{Violet}{Mechanistic toxicology}}}}\\[0.6mm]
\href{https://doi.org/10.1016/j.jes.2025.12.074}{Ding} & PM$_{2.5}$ promotes vascular calcification via miR-27b-5p/HES1 signalling & \textit{J Environ Sci (China)} & In vitro + organoid + murine & PM$_{2.5}$ & China\\
\href{https://doi.org/10.1007/s00360-026-01698-5}{Maraschi} & Osmoregulatory disruption and metal bioaccumulation from settleable atmospheric PM & \textit{J Comp Physiol B} & Controlled exposure (ecotox) & Settleable APM & Brazil\\
\href{https://doi.org/10.1016/j.jes.2025.11.018}{Wu} & Macrophage NLRP3 inflammasome in real-ambient PM-induced hepatic glucose dysregulation & \textit{J Environ Sci (China)} & Real-ambient murine + co-culture & Real-ambient PM & China\\
\addlinespace[0.9mm]
\multicolumn{6}{@{}l@{}}{\cellcolor{Mist}\textsf{\textbf{\textcolor{Clay}{Occupational \& indoor}}}}\\[0.6mm]
\href{https://doi.org/10.3389/fmedt.2026.1796828}{Chen} & Nested suction instrument to cut intraoperative surgical-smoke PM$_{2.5}$ & \textit{Front Med Technol} & Bench / device experiment & PM$_{2.5}$ (surgical smoke) & China\\
\href{https://doi.org/10.1093/annweh/wxag055}{Cole} & Determinants of respirable crystalline silica exposure in Australian tunnelling & \textit{Ann Work Expo Health} & Observational exposure study & Respirable dust / RCS & Australia\\
\href{https://doi.org/10.1016/j.jhazmat.2026.143081}{Ding} & Stage-resolved hair Zn-isotope dynamics as a firefighting combustion-exposure tracer & \textit{J Hazard Mater} & Repeated-measures biomarker & Combustion-derived PM & China\\
\href{https://doi.org/10.3390/jox16040119}{Ramsden} & Xenobiotic hazards in aircraft cabin air & \textit{J Xenobiot} & Review + risk economics & Ultrafine particles & Global\\
\addlinespace[0.9mm]
\multicolumn{6}{@{}l@{}}{\cellcolor{Mist}\textsf{\textbf{\textcolor{Amber}{Sensing, forecasting \& instrumentation}}}}\\[0.6mm]
\href{https://doi.org/10.1016/j.jenvman.2026.130515}{Shi} & Adaptive spatiotemporal graph transformer for multi-site multi-step PM$_{2.5}$ forecasting & \textit{J Environ Manage} & Deep-learning model & PM$_{2.5}$ & China / India\\
\href{https://doi.org/10.3390/bios16070380}{Tsuchiya} & Seasonal and indoor-outdoor characteristics of fluorescent aerosol particles in offices & \textit{Biosensors} & Real-time sensor campaign & Bioaerosol / coarse & Japan\\
\addlinespace[0.9mm]
\multicolumn{6}{@{}l@{}}{\cellcolor{Mist}\textsf{\textbf{\textcolor{Amber}{Exposure assessment \& modelling}}}}\\[0.6mm]
\href{https://doi.org/10.1029/2025GH001621}{Solorzano-Araque} & Multiscale exposure modelling: LOTOS-EUROS + mobility surveys + particle toxicity & \textit{GeoHealth} & CTM + mobility integration & PM$_{2.5}$ / PM$_{10}$ & Colombia\\
\addlinespace[0.9mm]
\multicolumn{6}{@{}l@{}}{\cellcolor{Mist}\textsf{\textbf{\textcolor{Coral}{Burden, policy \& mitigation}}}}\\[0.6mm]
\href{https://doi.org/10.1002/hsr2.72886}{Karamzad} & Ischaemic heart disease burden in MENA 1990-2021; PM attributable fraction 28.2\% & \textit{Health Sci Rep} & GBD secondary analysis & Ambient PM & MENA\\
\href{https://doi.org/10.1371/journal.pone.0347989}{Rees} & HEPA filtration to reduce respiratory infection in care homes & \textit{PLoS One} & RCT process evaluation & Filtered PM & UK\\
\href{https://doi.org/10.1016/j.jes.2026.01.009}{Shan} & Policy-driven emission reductions and post-pandemic rebound, high-resolution inventories & \textit{J Environ Sci (China)} & Emission inventory & PM emissions & China\\
\addlinespace[0.9mm]
\multicolumn{6}{@{}l@{}}{\cellcolor{Mist}\textsf{\textbf{\textcolor{Slate}{Other clinical endpoints}}}}\\[0.6mm]
\href{https://doi.org/10.3390/dj14070440}{Lupita} & PM$_{2.5}$, preventive-care capacity and caries/edentulism burden across the EU & \textit{Dent J} & Ecological panel & PM$_{2.5}$ & EU-27\\
\href{https://doi.org/10.1016/j.jdin.2026.06.008}{Park} & Long-term PM exposure and risk of androgenetic alopecia & \textit{JAAD Int} & Nationwide cohort & PM$_{2.5}$ / PM$_{10}$ & South Korea\\
\href{https://doi.org/10.3171/2026.3.PEDS2615}{Sun} & Environmental exposure and area deprivation vs. paediatric hydrocephalus outcomes & \textit{J Neurosurg Pediatr} & Registry cohort + geospatial & PM$_{2.5}$ (satellite-derived) & USA\\
\addlinespace[0.9mm]
\bottomrule
\end{longtable}}


\vspace{4mm}
\begin{keybox}
{\sffamily\bfseries\footnotesize\color{Deep}Provenance and method}\\[1.2mm]
{\fontsize{7.2}{8.8}\selectfont
Entry window \textbf{2026-08-25 to 2026-08-25}, a single complete day of deposits.
\textbf{Sources.} PubMed E-utilities on the \texttt{[Date - Entry]} field, run as two
separate queries (health axis 9 hits, instrumentation axis 0); the PubMed connector
returned 13 on the health axis and 0 on the instrumentation axis. Europe PMC
\texttt{CREATION\_DATE} 13 hits. Crossref by-ISSN across the 24 tracked
non-PubMed-indexed titles, 43 deposits after front-matter suppression. OpenAlex returned
empty for the 23rd consecutive run --- the date-filter leg has been broken since July and
is not being repaired. arXiv contributed nothing new.
\textbf{Screening.} 74 raw records, 65 unique after PMID/DOI/title collapse, 2 already
present in \texttt{seen.json}, \textbf{14 retained in scope and 51 rejected} with reason
logged. Rejections were dominated by non-particle atmospheric chemistry (ozone, VOC,
ammonia, gravity waves), non-inhalation ecotoxicology, and eleven Elsevier retraction
notices.
\textbf{Preprints.} Five records are unreviewed preprints (medRxiv, Research Square).
Tier reflects study design only and carries no statement about review status.
\textbf{Effect estimates.} All four interval-bearing estimates come from one preprint and
mix an odds ratio with three conditional exceedance ratios; they are plotted on a common
axis for precision and direction and are not mutually comparable. Estimates from
Pathak et al.\ are mean differences on an unbounded scale and are reported in the digest
text rather than plotted.
\textbf{Labels.} Subtopic, design, particle-metric and geography are assigned from title,
abstract and keywords; assignment is single-label by dominant contribution and is a
judgement, not a controlled vocabulary. All DOI links were resolved before build.\par}
\end{keybox}

\end{document}
